\documentclass[../../main.tex]{subfiles}

\begin{document}

\section{集合的并与交}\label{section:ch00-集合并交}

对于集合 $A$\footnote{《牛津英语字典》花了几百页的篇幅给出“集合”一词的定义.}, 元素 $x$ 是 $A$ 的成员关系记为 $x\in A$, 而 $x$ 不是 $A$ 的成员关系记为 $x\notin A$. 我们常说 $A$ 的一个成员属于 $A$ 且称 $A$ 的成员是 $A$ 中的一个点. 通常集合用花括号表示, 因此 $\{x\mid\text{关于 }x\text{ 的陈述}\}$ 是使得关于 $x$ 的陈述成立的所有元素 $x$ 的集合.

若两个集合有相同的成员, 我们说它们相同. 令 $A$ 和 $B$ 为集合. 若 $A$ 的每个成员也是 $B$ 的成员, 我们称 $A$ 为 $B$ 的子集, 记之为 $A\subseteq B$, 也说 $A$ 包含于 $B$ 或 $B$ 包含 $A$. $B$ 的子集 $A$ 称为 $B$ 的真子集, 若 $A\neq B$. $A$ 和 $B$ 的并, 记为 $A\cup B$, 是所有或者属于 $A$ 或者属于 $B$ 的点的集合, 即
\begin{align*}
A\cup B=\{x\mid x\in A\text{ 或 }x\in B\}.
\end{align*}
这里“或”这个词在非互斥的意义下使用, 因此同时属于 $A$ 和 $B$ 的点属于 $A\cup B$. $A$ 和 $B$ 的交, 记为 $A\cap B$, 是所有同时属于 $A$ 和 $B$ 的点的集合, 即
\begin{align*}
A\cap B=\{x\mid x\in A\text{ 且 }x\in B\}.
\end{align*}
$A$ 在 $B$ 中的补, 记为 $B\setminus A$, 是 $B$ 中那些不在 $A$ 中的点的集合, 即
\begin{align*}
B\setminus A=\{x\in B\mid x\notin A\}.
\end{align*}
若在特别的讨论中所有的集合是参考集 $X$ 的子集, 我们常简单地称 $X\setminus A$ 为 $A$ 的补.

没有任何成员的集合称为空集, 记为 $\varnothing$. 不等于空集的集合称为非空的. 我们称只有一个成员的集合为单点集. 给定集合 $X$, $X$ 的所有子集的集合记为 $\mathcal{P}(X)$ 或 $2^{X}$, 称之为 $X$ 的幂集.

为了避免考虑集合的集合时可能产生混淆, 我们常用词“族”或“簇”作为“集”的同义词. 令 $\mathcal{F}$ 为集族. $\mathcal{F}$ 的并, 记为 $\bigcup_{F\in\mathcal{F}}F$, 定义为属于 $\mathcal{F}$ 中的至少一个集合的点的集合. $\mathcal{F}$ 的交, 记为 $\bigcap_{F\in\mathcal{F}}F$, 定义为属于 $\mathcal{F}$ 中的每个集合的点的集合. 若集族 $\mathcal{F}$ 中的任何两个集合的交是空的, 集族 $\mathcal{F}$ 称为是不交的. 对于集族 $\mathcal{F}$, 通过检验集合的包含关系可得到以下等式.

\begin{proposition}[De Morgan 等式]\label{proposition:ch00-demorgan}
对于参考集 $X$ 与集族 $\mathcal{F}$, 有
\begin{align*}
X\setminus\left[\bigcup_{F\in\mathcal{F}}F\right]=\bigcap_{F\in\mathcal{F}}[X\setminus F],\qquad
X\setminus\left[\bigcap_{F\in\mathcal{F}}F\right]=\bigcup_{F\in\mathcal{F}}[X\setminus F].
\end{align*}
即并的补是补的交, 且交的补是补的并.
\end{proposition}

对于集合 $\Lambda$, 假定对每个 $\lambda\in\Lambda$, 存在已定义的 $E_{\lambda}$. 令 $\mathcal{F}$ 为集族 $\{E_{\lambda}\mid\lambda\in\Lambda\}$. 我们写作 $\mathcal{F}=\{E_{\lambda}\}_{\lambda\in\Lambda}$ 且称之为 $\mathcal{F}$ 的用指标集（或参数集）$\Lambda$ 标记的指标（或参数化）.

\end{document}